\scanpage{27}
% L1–2
\uline{Cube}

\textcircled{A} \uline{Pt de vue géométrie des} \unsure{$\mathbb R$|P} \uline{convexes} (sur $\mathbb R$)

% L3–4
$\circledast$ $n$-Cube \unsure{standart|standard} (réel) :
\[
\begin{aligned}
 C_n&\subset\mathbb R^n\\[-2pt]
 \Vert\quad&\\[-4pt]
 \{x=(x_i)_{1\leq i\leq n}\in\mathbb R^n&\mid |x_i|\leq1\}
\end{aligned}
\]

% L5–6
$\circledast$ $n$-Cube \inserted{$C$} d'un \deleted{vectoriel} \inserted{espace affine} réel \inserted{$E$} (de dim finie $n$) : partie
telle que il existe isom \inserted{aff} $\mathbb R^n\overset{u}{\unsure{-|\longrightarrow}}E$ avec $u(C_n)=C$.

% L7–11
$\circledast$ Catégorie des $n$-cubes \deleted{de dim $n$} \inserted{convexes} (des \deleted{vectoriels} \inserted{$\mathbb R$-espaces affines} de dim $n$) :
\deleted{\unsure{est}} c'est une groupoïde connexe avec objet
marqué \deleted{\unsure{$C_n$}} $(\mathbb R^n,C_n)$, donc équivaut à la
catégorie des torseurs sous
\[
 \deleted{\unsure{B_n}} W_n=\mathrm{Aut}_{\deleted{\unsure{\mathrm{vect}.}}\mathrm{aff}}(\mathbb R^n,C_n)
\]

% L12–14
% margin note written sideways
\margin{expliciter en termes de $B\overset{p}{\longrightarrow}I=B/R$\par
$E=\mathrm{Hom}_{\{\pm1\}}(B,\mathbb R)$\par
$C=\{\lambda\in\mathrm{Hom}_{\{\pm1\}}(B,\mathbb R)\mid |\lambda(b)|\leq1\quad\forall b\in B\}$}
\uline{Prop.} Soit $B_n=\{\pm e_i\mid 1\leq i\leq n\}\subset\mathbb R^n$. Alors
\[
 \mathrm{Aut}_{\mathrm{aff}}(\mathbb R^n,C_n)=\mathrm{Aut}_{\mathrm{vect}}(\mathbb R^{\deleted{\unsure{n}}n},B_n)
 \xrightarrow{\sim}\mathrm{Aut}(B_n,\deleted{\unsure{R_n}}\unsure{\sigma_n|R_n})
\]
\margin{\unsure{rel|relat} \deleted{\illegible{1}}}
où $\sigma_n$ est l'involution $b\mapsto-b$ de $B_n$.

% L15–18
\deleted{\unsure{Dém}} Dém. laissée comme exercice. [ $\supset$ est
évident, mais pas tout à fait $\subset$ ; on peut noter
que $B_n$ est l'ens \unsure{des|ses} \uline{pts extrémaux} du
convexe $C_n$\ldots{} ]

% L19–22
\uline{Corollaire} Catégorie des $n$-cubes convexes équivaut
à \deleted{\illegible{1}} la catégorie \deleted{des} \inserted{épures de} $n$-cubes \deleted{combinatoires}.

\uline{Corollaire} Cat. des $n$-cubes orientés $\cong$ cat. des épures
de $n$-cubes \deleted{combinatoires} orientés.
% archive sheet number: 26
